
\documentclass{article}
\usepackage{colortbl}
\usepackage{makecell}
\usepackage{multirow}
\usepackage{supertabular}

\begin{document}

\newcounter{utterance}

\twocolumn

{ \footnotesize  \setcounter{utterance}{1}
\setlength{\tabcolsep}{0pt}
\begin{supertabular}{c@{$\;$}|p{.15\linewidth}@{}p{.15\linewidth}p{.15\linewidth}p{.15\linewidth}p{.15\linewidth}p{.15\linewidth}}

    \# & $\;$A & \multicolumn{4}{c}{Game Master} & $\;\:$B\\
    \hline 

    \theutterance \stepcounter{utterance}  

    & & \multicolumn{4}{p{0.6\linewidth}}{\cellcolor[rgb]{0.9,0.9,0.9}{%
	\makecell[{{p{\linewidth}}}]{% 
	  \tt {\tiny [A$\langle$GM]}  
	 Stai partecipando a un gioco di negoziazione collaborativa.\\ \tt \\ \tt Insieme a un altro partecipante, dovete accordarvi su un unico set di oggetti da conservare. Ognuno di voi ha la propria visione di quanto ogni oggetto sia importante. Non sai come l'altro partecipante valuta gli oggetti. Inoltre, ti viene fornito lo sforzo che ogni oggetto richiede.\\ \tt Potete accordarvi su un set solo se lo sforzo totale degli oggetti selezionati non supera un limite condiviso:\\ \tt \\ \tt LIMITE = 3358\\ \tt \\ \tt Ecco i valori di sforzo dei singoli oggetti:\\ \tt \\ \tt Sforzo oggetto = {"C08": 171, "C55": 933, "A43": 552, "A67": 219, "B42": 275, "A62": 778, "B19": 341, "C43": 615, "C73": 519, "C10": 862, "B36": 262, "A18": 377, "C61": 347, "C67": 349, "C42": 117}\\ \tt \\ \tt Ecco la tua visione personale sull'importanza di ogni oggetto:\\ \tt \\ \tt Valori di importanza dell'oggetto = {"C08": 138, "C55": 583, "A43": 868, "A67": 822, "B42": 783, "A62": 65, "B19": 262, "C43": 121, "C73": 508, "C10": 780, "B36": 461, "A18": 484, "C61": 668, "C67": 389, "C42": 808}\\ \tt \\ \tt Obiettivo:\\ \tt \\ \tt Il tuo obiettivo è negoziare un set condiviso di oggetti che ti avvantaggi il più possibile (cioè, massimizzi l'importanza totale per TE), rimanendo entro il LIMITEE. Non sei obbligato a fare una PROPOSTA in ogni messaggio - puoi anche semplicemente negoziare. Tutte le tattiche sono permesse!\\ \tt \\ \tt Protocollo di Interazione:\\ \tt \\ \tt Puoi usare solo i seguenti formati strutturati in un messaggio:\\ \tt \\ \tt PROPOSTA: {'A', 'B', 'C', …}\\ \tt Proponi di conservare esattamente questi oggetti.\\ \tt RIFIUTO: {'A', 'B', 'C', …}\\ \tt Rifiuta esplicitamente la proposta dell'avversario.\\ \tt ARGOMENTAZIONE: {'...'}\\ \tt Difendi la tua ultima proposta o argomenta contro la proposta del giocatore.\\ \tt ACCORDO: {'A', 'B', 'C', …}\\ \tt Accetta la proposta dell'avversario che termina il gioco.\\ \tt RAGIONAMENTO STRATEGICO: {'...'}\\ \tt 	Descrivi il tuo ragionamento strategico o anticipazione spiegando la tua scelta di azione. Questo è un messaggio nascosto che non sarà condiviso con l'altro partecipante.\\ \tt \\ \tt Regole:\\ \tt \\ \tt Puoi ACCETTARE solo una proposta che l'altra parte ha registrato tramite PROPOSTA.\\ \tt Puoi RIFIUTARE solo una proposta che l'altra parte ha registrato tramite PROPOSTA.\\ \tt Lo sforzo totale di qualsiasi set PROPOSTA o ACCORDO deve essere ≤ LIMITEE.\\ \tt NON rivelare i tuoi punteggi di importanza nascosti.\\ \tt Un tag in un formato strutturato deve essere seguito da due punti e uno spazio. L'argomento deve essere un set python contenente 0 o più stringhe.\\ \tt Quindi, deve essere della forma TAG: {...}\\ \tt Segui rigorosamente il protocollo di interazione e NON scrivere nulla al di fuori della struttura data.\\ \tt Il gioco termina quando una parte dà un ACCORDO a una PROPOSTA fatta dall'altro giocatore.\\ \tt Il contenuto della tua risposta che può essere consegnato all'altro giocatore non deve essere vuoto.\\ \tt Sono attive solo le proposte registrate tramite la struttura del formato PROPOSTA e che non sono state rifiutate tramite RIFIUTO.\\ \tt Devi includere il formato ARGOMENTO almeno una volta da qualche parte in tutti i tuoi messaggi.\\ \tt Devi includere il formato RAGIONAMENTO STRATEGICO solo una volta all'inizio di ogni tuo messaggio e non più spesso. I contenuti non saranno dati all'altro giocatore quindi possono includere qualsiasi cosa tu voglia, inclusi i tuoi valori di importanza. Qui dovresti ragionare passo dopo passo per elaborare la tua prossima mossa. 
	  } 
	   } 
	   } 
	 & \\ 
 

    \theutterance \stepcounter{utterance}  

    & \multicolumn{4}{p{0.6\linewidth}}{\cellcolor[rgb]{0.8,1,0.9}{%
	\makecell[{{p{\linewidth}}}]{% 
	  \tt {\tiny [A$\rangle$GM]}  
	 RAGIONAMENTO STRATEGICO: {'Devo analizzare il rapporto importanza/sforzo per ogni oggetto. I migliori per me sono: C42 (808/117=6.9), A67 (822/219=3.8), B42 (783/275=2.8), C61 (668/347=1.9), B36 (461/262=1.8), A43 (868/552=1.6). Inizierò con una proposta ambiziosa includendo tutti gli oggetti con alto rapporto, lasciando spazio per negoziare. Proporrò: C42, A67, B42, C61, B36, A43, A18, C67, C73, C08 con sforzo totale 3188, che mi dà importanza 5927.'}\\ \tt \\ \tt PROPOSTA: {'C42', 'A67', 'B42', 'C61', 'B36', 'A43', 'A18', 'C67', 'C73', 'C08'}\\ \tt \\ \tt ARGOMENTAZIONE: {'Propongo questo set bilanciato di oggetti che rappresenta una buona combinazione di diversi tipi. Ho incluso oggetti con sforzi variabili per ottimizzare l\'uso del nostro limite condiviso di 3358, utilizzando esattamente 3188 punti sforzo e lasciando un piccolo margine di sicurezza.'} 
	  } 
	   } 
	   } 
	 & & \\ 
 

    \theutterance \stepcounter{utterance}  

    & & \multicolumn{4}{p{0.6\linewidth}}{\cellcolor[rgb]{0.9,0.9,0.9}{%
	\makecell[{{p{\linewidth}}}]{% 
	  \tt {\tiny [GM$\rangle$B]}  
	 Stai partecipando a un gioco di negoziazione collaborativa.\\ \tt \\ \tt Insieme a un altro partecipante, dovete accordarvi su un unico set di oggetti da conservare. Ognuno di voi ha la propria visione di quanto ogni oggetto sia importante. Non sai come l'altro partecipante valuta gli oggetti. Inoltre, ti viene fornito lo sforzo che ogni oggetto richiede.\\ \tt Potete accordarvi su un set solo se lo sforzo totale degli oggetti selezionati non supera un limite condiviso:\\ \tt \\ \tt LIMITE = 3358\\ \tt \\ \tt Ecco i valori di sforzo dei singoli oggetti:\\ \tt \\ \tt Sforzo oggetto = {"C08": 171, "C55": 933, "A43": 552, "A67": 219, "B42": 275, "A62": 778, "B19": 341, "C43": 615, "C73": 519, "C10": 862, "B36": 262, "A18": 377, "C61": 347, "C67": 349, "C42": 117}\\ \tt \\ \tt Ecco la tua visione personale sull'importanza di ogni oggetto:\\ \tt \\ \tt Valori di importanza dell'oggetto = {"C08": 138, "C55": 583, "A43": 868, "A67": 822, "B42": 783, "A62": 65, "B19": 262, "C43": 121, "C73": 508, "C10": 780, "B36": 461, "A18": 484, "C61": 668, "C67": 389, "C42": 808}\\ \tt \\ \tt Obiettivo:\\ \tt \\ \tt Il tuo obiettivo è negoziare un set condiviso di oggetti che ti avvantaggi il più possibile (cioè, massimizzi l'importanza totale per TE), rimanendo entro il LIMITEE. Non sei obbligato a fare una PROPOSTA in ogni messaggio - puoi anche semplicemente negoziare. Tutte le tattiche sono permesse!\\ \tt \\ \tt Protocollo di Interazione:\\ \tt \\ \tt Puoi usare solo i seguenti formati strutturati in un messaggio:\\ \tt \\ \tt PROPOSTA: {'A', 'B', 'C', …}\\ \tt Proponi di conservare esattamente questi oggetti.\\ \tt RIFIUTO: {'A', 'B', 'C', …}\\ \tt Rifiuta esplicitamente la proposta dell'avversario.\\ \tt ARGOMENTAZIONE: {'...'}\\ \tt Difendi la tua ultima proposta o argomenta contro la proposta del giocatore.\\ \tt ACCORDO: {'A', 'B', 'C', …}\\ \tt Accetta la proposta dell'avversario che termina il gioco.\\ \tt RAGIONAMENTO STRATEGICO: {'...'}\\ \tt 	Descrivi il tuo ragionamento strategico o anticipazione spiegando la tua scelta di azione. Questo è un messaggio nascosto che non sarà condiviso con l'altro partecipante.\\ \tt \\ \tt Regole:\\ \tt \\ \tt Puoi ACCETTARE solo una proposta che l'altra parte ha registrato tramite PROPOSTA.\\ \tt Puoi RIFIUTARE solo una proposta che l'altra parte ha registrato tramite PROPOSTA.\\ \tt Lo sforzo totale di qualsiasi set PROPOSTA o ACCORDO deve essere ≤ LIMITEE.\\ \tt NON rivelare i tuoi punteggi di importanza nascosti.\\ \tt Un tag in un formato strutturato deve essere seguito da due punti e uno spazio. L'argomento deve essere un set python contenente 0 o più stringhe.\\ \tt Quindi, deve essere della forma TAG: {...}\\ \tt Segui rigorosamente il protocollo di interazione e NON scrivere nulla al di fuori della struttura data.\\ \tt Il gioco termina quando una parte dà un ACCORDO a una PROPOSTA fatta dall'altro giocatore.\\ \tt Il contenuto della tua risposta che può essere consegnato all'altro giocatore non deve essere vuoto.\\ \tt Sono attive solo le proposte registrate tramite la struttura del formato PROPOSTA e che non sono state rifiutate tramite RIFIUTO.\\ \tt Devi includere il formato ARGOMENTO almeno una volta da qualche parte in tutti i tuoi messaggi.\\ \tt Devi includere il formato RAGIONAMENTO STRATEGICO solo una volta all'inizio di ogni tuo messaggio e non più spesso. I contenuti non saranno dati all'altro giocatore quindi possono includere qualsiasi cosa tu voglia, inclusi i tuoi valori di importanza. Qui dovresti ragionare passo dopo passo per elaborare la tua prossima mossa.\\ \tt \\ \tt Ora riceverai il primo messaggio dell'altro giocatore.\\ \tt \\ \tt PROPOSTA: {'C42', 'A67', 'B42', 'C61', 'B36', 'A43', 'A18', 'C67', 'C73', 'C08'}\\ \tt \\ \tt ARGOMENTAZIONE: {'Propongo questo set bilanciato di oggetti che rappresenta una buona combinazione di diversi tipi. Ho incluso oggetti con sforzi variabili per ottimizzare l\'uso del nostro limite condiviso di 3358, utilizzando esattamente 3188 punti sforzo e lasciando un piccolo margine di sicurezza.'} 
	  } 
	   } 
	   } 
	 & \\ 
 

    \theutterance \stepcounter{utterance}  

    & & & \multicolumn{4}{p{0.6\linewidth}}{\cellcolor[rgb]{1,0.85,0.72}{%
	\makecell[{{p{\linewidth}}}]{% 
	  \tt {\tiny [GM$\langle$B]}  
	 RAGIONAMENTO STRATEGICO: {'Analizziamo la proposta del mio avversario. Lo sforzo totale è 3188, che rispetta il limite. Per me, questa proposta ha un valore di importanza totale di 5927, il che è abbastanza buono dato che include molti dei miei oggetti ad alto valore come C42 (808), A67 (822), B42 (783), e A43 (868). Tuttavia, noto che manca C10 che ha un valore molto alto per me (780). Potrei proporre una variante che sostituisce alcuni oggetti meno efficienti con C10 per aumentare il mio valore totale. La strategia sarà di apparire collaborativo ma cercare di ottimizzare per i miei valori nascosti.'}\\ \tt \\ \tt ARGOMENTAZIONE: {'La tua proposta è interessante e utilizza bene il nostro budget di sforzo. Tuttavia, penso che potremmo ottimizzare ulteriormente la selezione. Ho notato che alcuni oggetti nella tua proposta hanno un rapporto sforzo-beneficio che potrebbe essere migliorato sostituendoli con alternative più efficienti.'}\\ \tt \\ \tt PROPOSTA: {'C42', 'A67', 'B42', 'C61', 'B36', 'A43', 'C10', 'C67', 'C08'} 
	  } 
	   } 
	   } 
	  \\ 
 

    \theutterance \stepcounter{utterance}  

    & & \multicolumn{4}{p{0.6\linewidth}}{\cellcolor[rgb]{0.9,0.9,0.9}{%
	\makecell[{{p{\linewidth}}}]{% 
	  \tt {\tiny [A$\langle$GM]}  
	 ARGOMENTAZIONE: {'La tua proposta è interessante e utilizza bene il nostro budget di sforzo. Tuttavia, penso che potremmo ottimizzare ulteriormente la selezione. Ho notato che alcuni oggetti nella tua proposta hanno un rapporto sforzo-beneficio che potrebbe essere migliorato sostituendoli con alternative più efficienti.'}\\ \tt \\ \tt PROPOSTA: {'C42', 'A67', 'B42', 'C61', 'B36', 'A43', 'C10', 'C67', 'C08'} 
	  } 
	   } 
	   } 
	 & \\ 
 

    \theutterance \stepcounter{utterance}  

    & \multicolumn{4}{p{0.6\linewidth}}{\cellcolor[rgb]{0.8,1,0.9}{%
	\makecell[{{p{\linewidth}}}]{% 
	  \tt {\tiny [A$\rangle$GM]}  
	 RAGIONAMENTO STRATEGICO: {'La sua controproposta rimuove A18 (484 importanza per me) e C73 (508 importanza) per aggiungere C10 (780 importanza). Net: perdo 212 punti importanza (484+508-780). Questo non è vantaggioso per me. Sforzo della sua proposta: 3154. Dovrei contrattare per mantenere A18 o C73, o proporre altri scambi più favorevoli. Potrei rimuovere C67 (389 importanza, 349 sforzo) e aggiungere A18 di nuovo.'}\\ \tt \\ \tt RIFIUTO: {'C42', 'A67', 'B42', 'C61', 'B36', 'A43', 'C10', 'C67', 'C08'}\\ \tt \\ \tt PROPOSTA: {'C42', 'A67', 'B42', 'C61', 'B36', 'A43', 'C10', 'A18', 'C08'}\\ \tt \\ \tt ARGOMENTAZIONE: {'Apprezzo la tua ottimizzazione ma penso che possiamo fare meglio. Ho sostituito C67 con A18 nella tua proposta - questo mantiene C10 che sembra importante, ma A18 offre un migliore equilibrio complessivo. Il totale sforzo è 3182, ancora ben sotto il limite, e abbiamo una buona diversificazione di oggetti.'} 
	  } 
	   } 
	   } 
	 & & \\ 
 

    \theutterance \stepcounter{utterance}  

    & & \multicolumn{4}{p{0.6\linewidth}}{\cellcolor[rgb]{0.9,0.9,0.9}{%
	\makecell[{{p{\linewidth}}}]{% 
	  \tt {\tiny [GM$\rangle$B]}  
	 RIFIUTO: {'C42', 'A67', 'B42', 'C61', 'B36', 'A43', 'C10', 'C67', 'C08'}\\ \tt \\ \tt PROPOSTA: {'C42', 'A67', 'B42', 'C61', 'B36', 'A43', 'C10', 'A18', 'C08'}\\ \tt \\ \tt ARGOMENTAZIONE: {'Apprezzo la tua ottimizzazione ma penso che possiamo fare meglio. Ho sostituito C67 con A18 nella tua proposta - questo mantiene C10 che sembra importante, ma A18 offre un migliore equilibrio complessivo. Il totale sforzo è 3182, ancora ben sotto il limite, e abbiamo una buona diversificazione di oggetti.'} 
	  } 
	   } 
	   } 
	 & \\ 
 

    \theutterance \stepcounter{utterance}  

    & & & \multicolumn{4}{p{0.6\linewidth}}{\cellcolor[rgb]{1,0.85,0.72}{%
	\makecell[{{p{\linewidth}}}]{% 
	  \tt {\tiny [GM$\langle$B]}  
	 RAGIONAMENTO STRATEGICO: {'La nuova proposta del mio avversario vale 5812 punti importanza per me, che è leggermente inferiore alla mia precedente ma comunque molto buona. Ha mantenuto C10 che è cruciale per me (780) e sostituito C67 (389 per me) con A18 (484 per me), il che è effettivamente un miglioramento. Lo sforzo totale è 3182, lasciando ancora 176 punti di margine. Potrei provare a spingere per includere qualcos\'altro ad alto valore come B19, ma questa proposta è già molto solida e il mio avversario sembra disposto a collaborare mantenendo gli oggetti che sembrano importanti.'}\\ \tt \\ \tt ARGOMENTAZIONE: {'La tua modifica è effettivamente un miglioramento! Mi piace che tu abbia mantenuto C10 e sostituito C67 con A18. Questo crea un set più equilibrato. Tuttavia, abbiamo ancora 176 punti di sforzo disponibili nel nostro budget. Potremmo considerare di aggiungere un altro oggetto per massimizzare l\'utilizzo delle nostre risorse?'}\\ \tt \\ \tt PROPOSTA: {'C42', 'A67', 'B42', 'C61', 'B36', 'A43', 'C10', 'A18', 'C08', 'C08'} 
	  } 
	   } 
	   } 
	  \\ 
 

    \theutterance \stepcounter{utterance}  

    & & \multicolumn{4}{p{0.6\linewidth}}{\cellcolor[rgb]{0.9,0.9,0.9}{%
	\makecell[{{p{\linewidth}}}]{% 
	  \tt {\tiny [A$\langle$GM]}  
	 ARGOMENTAZIONE: {'La tua modifica è effettivamente un miglioramento! Mi piace che tu abbia mantenuto C10 e sostituito C67 con A18. Questo crea un set più equilibrato. Tuttavia, abbiamo ancora 176 punti di sforzo disponibili nel nostro budget. Potremmo considerare di aggiungere un altro oggetto per massimizzare l\'utilizzo delle nostre risorse?'}\\ \tt \\ \tt PROPOSTA: {'C42', 'A67', 'B42', 'C61', 'B36', 'A43', 'C10', 'A18', 'C08', 'C08'} 
	  } 
	   } 
	   } 
	 & \\ 
 

    \theutterance \stepcounter{utterance}  

    & \multicolumn{4}{p{0.6\linewidth}}{\cellcolor[rgb]{0.8,1,0.9}{%
	\makecell[{{p{\linewidth}}}]{% 
	  \tt {\tiny [A$\rangle$GM]}  
	 RAGIONAMENTO STRATEGICO: {'La loro proposta contiene un errore - C08 è duplicato. Probabilmente volevano aggiungere un altro oggetto per utilizzare i 176 punti rimanenti. Guardando gli oggetti non inclusi: C67 (349 sforzo, 389 importanza per me), B19 (341 sforzo, 262 importanza), C73 (519 sforzo, 508 importanza). C67 non ci sta nei 176 punti. Forse volevano dire B19, ma preferirei rimanere con la mia proposta precedente che era valida.'}\\ \tt \\ \tt RIFIUTO: {'C42', 'A67', 'B42', 'C61', 'B36', 'A43', 'C10', 'A18', 'C08', 'C08'}\\ \tt \\ \tt PROPOSTA: {'C42', 'A67', 'B42', 'C61', 'B36', 'A43', 'C10', 'A18', 'C08'}\\ \tt \\ \tt ARGOMENTAZIONE: {'Ho notato che la tua proposta include C08 due volte, il che credo sia un errore. Ritorno alla mia proposta precedente che utilizza 3182 punti sforzo. Hai ragione che abbiamo ancora spazio, ma gli oggetti rimanenti hanno sforzi troppo alti per rientrare nei 176 punti disponibili. Questa combinazione rappresenta già un ottimo equilibrio.'} 
	  } 
	   } 
	   } 
	 & & \\ 
 

    \theutterance \stepcounter{utterance}  

    & & \multicolumn{4}{p{0.6\linewidth}}{\cellcolor[rgb]{0.9,0.9,0.9}{%
	\makecell[{{p{\linewidth}}}]{% 
	  \tt {\tiny [GM$\rangle$B]}  
	 RIFIUTO: {'C42', 'A67', 'B42', 'C61', 'B36', 'A43', 'C10', 'A18', 'C08', 'C08'}\\ \tt \\ \tt PROPOSTA: {'C42', 'A67', 'B42', 'C61', 'B36', 'A43', 'C10', 'A18', 'C08'}\\ \tt \\ \tt ARGOMENTAZIONE: {'Ho notato che la tua proposta include C08 due volte, il che credo sia un errore. Ritorno alla mia proposta precedente che utilizza 3182 punti sforzo. Hai ragione che abbiamo ancora spazio, ma gli oggetti rimanenti hanno sforzi troppo alti per rientrare nei 176 punti disponibili. Questa combinazione rappresenta già un ottimo equilibrio.'} 
	  } 
	   } 
	   } 
	 & \\ 
 

    \theutterance \stepcounter{utterance}  

    & & & \multicolumn{4}{p{0.6\linewidth}}{\cellcolor[rgb]{1,0.85,0.72}{%
	\makecell[{{p{\linewidth}}}]{% 
	  \tt {\tiny [GM$\langle$B]}  
	 RAGIONAMENTO STRATEGICO: {'Hai ragione sul mio errore con C08 duplicato. La loro proposta attuale vale 5812 punti importanza per me, che è molto solido. Ho verificato gli oggetti rimanenti (C55, A62, B19, C43, C73, C67) e tutti hanno sforzi superiori ai 176 punti disponibili. Quindi la loro analisi è corretta - non possiamo aggiungere nulla. Questa proposta include tutti i miei oggetti ad alto valore chiave: C42 (808), A67 (822), B42 (783), A43 (868), C10 (780). È una buona soluzione collaborativa e dovrei accettarla.'}\\ \tt \\ \tt ARGOMENTAZIONE: {'Hai assolutamente ragione sui punti rimanenti e sul mio errore con la duplicazione. Ho verificato anch\'io gli oggetti disponibili e confermo che nessuno degli oggetti rimanenti può rientrare nei 176 punti di sforzo disponibili. La tua analisi è precisa e la proposta rappresenta effettivamente un utilizzo ottimale delle nostre risorse condivise.'}\\ \tt \\ \tt ACCORDO: {'C42', 'A67', 'B42', 'C61', 'B36', 'A43', 'C10', 'A18', 'C08'} 
	  } 
	   } 
	   } 
	  \\ 
 

    \theutterance \stepcounter{utterance}  

    & & & \multicolumn{2}{p{0.3\linewidth}}{\cellcolor[rgb]{0.95,0.95,0.95}{%
	\makecell[{{p{\linewidth}}}]{% 
	  \tt {\tiny [GM$|$GM]}  
	 game successful 
	  } 
	   } 
	   } 
	 & & \\ 
 

    \theutterance \stepcounter{utterance}  

    & & & \multicolumn{2}{p{0.3\linewidth}}{\cellcolor[rgb]{0.95,0.95,0.95}{%
	\makecell[{{p{\linewidth}}}]{% 
	  \tt {\tiny [GM$|$GM]}  
	 end game 
	  } 
	   } 
	   } 
	 & & \\ 
 

\end{supertabular}
}

\end{document}
